\documentclass{article}
\usepackage{amsmath,amssymb,amsthm,algorithm,algorithmic}
\usepackage{bm}
\newtheorem{theorem}{Theorem}
\newtheorem{lemma}{Lemma}
\newtheorem{assumption}{Assumption}
\newcommand{\norm}[1]{\left\lVert#1\right\rVert}
\begin{document}

\section*{Primal–Dual Hybrid Gradient (PDHG) Algorithm}

\section*{Mathematical Formulation}
We consider the convex optimisation problem  

\[
\min_{x\in\mathbb R^{n}} \; \underbrace{f(x)}_{\text{smooth convex}} \;+\; \underbrace{g(Kx)}_{\text{convex (possibly nonsmooth)}},
\tag{P}
\]

where  

\begin{itemize}
    \item $f:\mathbb R^{n}\to\mathbb R$ is convex and $L$‑smooth, i.e. 
          $\|\nabla f(x)-\nabla f(x')\|\le L\norm{x-x'}$ for all $x,x'$,
    \item $g:\mathbb R^{m}\to\mathbb R\cup\{+\infty\}$ is convex, proper and lower‑semicontinuous,
    \item $K\in\mathbb R^{m\times n}$ is a linear operator with operator norm $\|K\|$.
\end{itemize}

Introducing the Fenchel conjugate $g^{*}$ and the dual variable $y\in\mathbb R^{m}$, the saddle‑point formulation reads  

\[
\min_{x}\max_{y}\; \mathcal L(x,y)\;:=\; f(x)+\langle Kx,y\rangle-g^{*}(y).
\tag{1}
\]

The PDHG iterates are

\[
\boxed{
\begin{aligned}
y_{k+1} &= \operatorname{prox}_{\sigma g^{*}}\!\bigl(y_{k}+ \sigma K x_{k}\bigr), \tag{2a}\\[2mm]
x_{k+1} &= \operatorname{prox}_{\tau f}\!\bigl(x_{k}-\tau K^{\top} y_{k+1}\bigr), \tag{2b}
\end{aligned}}
\]

with step sizes $\tau>0$, $\sigma>0$ satisfying the **stability condition**

\[
\tau\sigma\|K\|^{2}<1.
\tag{3}
\]

When $g$ is smooth, the proximal step on the dual reduces to a simple gradient ascent:

\[
y_{k+1}= y_{k}+ \sigma\bigl(Kx_{k}-\nabla g^{*}(y_{k})\bigr).
\tag{4}
\]

\section*{Pseudocode}
\begin{algorithm}[H]
\caption{Primal–Dual Hybrid Gradient (PDHG)}
\begin{algorithmic}[1]
\STATE \textbf{Input:} $x_{0}\in\mathbb R^{n}$, $y_{0}\in\mathbb R^{m}$, stepsizes $\tau,\sigma>0$ satisfying (3)
\FOR{$k=0,1,2,\dots$}
    \STATE $\,\displaystyle y_{k+1}\gets \operatorname{prox}_{\sigma g^{*}}\!\bigl(y_{k}+ \sigma K x_{k}\bigr)$
    \STATE $\,\displaystyle x_{k+1}\gets \operatorname{prox}_{\tau f}\!\bigl(x_{k}-\tau K^{\top} y_{k+1}\bigr)$
\ENDFOR
\STATE \textbf{Output:} last iterate $(x_{K},y_{K})$ or averaged iterate $\bar x_{K}=\frac1K\sum_{k=1}^{K}x_{k}$.
\end{algorithmic}
\end{algorithm}

\section*{Dry Run Example}
Consider the scalar problem  

\[
\min_{w\in\mathbb R}\;(w-3)^{2},
\]

i.e. $f(w)=(w-3)^{2}$, $g\equiv0$, $K=1$.  
Since $g^{*}=0$, $\operatorname{prox}_{\sigma g^{*}}$ is the identity, thus $y_{k}=0$ for all $k$.  
Choose $\tau=\sigma=\eta=0.1$ (note that $\tau\sigma\|K\|^{2}=0.01<1$).

\[
\begin{array}{c|c|c|c}
k & x_{k} & \text{gradient } \nabla f(x_{k}) & f(x_{k})\\ \hline
0 & 0 & 2(0-3)=-6 & 9\\
1 & x_{1}=x_{0}-\tau\nabla f(x_{0})=0-0.1(-6)=0.6 & 2(0.6-3)=-4.8 & (0.6-3)^{2}=5.76\\
2 & x_{2}=0.6-0.1(-4.8)=1.08 & 2(1.08-3)=-3.84 & (1.08-3)^{2}=3.6864\\
3 & x_{3}=1.08-0.1(-3.84)=1.464 & 2(1.464-3)=-3.072 & (1.464-3)^{2}=2.3660
\end{array}
\]

The objective value drops from $9$ to $2.37$ after three iterations, illustrating the gradient‑type behaviour of PDHG when the dual term is absent.

\newpage
\section*{Assumptions}
\begin{assumption}[Convexity and Smoothness]\label{as:convex}
\begin{enumerate}
\item $f$ is convex and $L$‑smooth:
\[
f(x')\le f(x)+\langle\nabla f(x),x'-x\rangle+\frac{L}{2}\norm{x'-x}^{2},\quad\forall x,x'.
\tag{A1}
\]
\item $g$ is convex, proper, lower‑semicontinuous and its Fenchel conjugate $g^{*}$ is $\frac{1}{\mu}$‑strongly convex (i.e. $g$ is $\mu$‑smooth). When $g$ is nonsmooth we only require that the proximal operator $\operatorname{prox}_{\sigma g^{*}}$ is well defined.
\item The linear operator $K$ satisfies $\|K\|<\infty$.
\end{enumerate}
\end{assumption}

\begin{assumption}[Step‑size Condition]\label{as:step}
The step sizes $\tau,\sigma>0$ obey  

\[
\tau\sigma\|K\|^{2}\le \theta<1,
\tag{A2}
\]

for a prescribed constant $\theta\in(0,1)$.  
\end{assumption}

\section*{Main Convergence Theorem}
\begin{theorem}[Ergodic $O(1/T)$ convergence]\label{thm:ergodic}
Let $(x_{k},y_{k})_{k\ge0}$ be generated by \eqref{2a}–\eqref{2b} under Assumptions \ref{as:convex}–\ref{as:step}.  
Define the ergodic averages  

\[
\bar x_{T}:=\frac{1}{T}\sum_{k=0}^{T-1}x_{k},\qquad 
\bar y_{T}:=\frac{1}{T}\sum_{k=0}^{T-1}y_{k}.
\]

Then for any saddle point $(x^{\star},y^{\star})$ of $\mathcal L$,

\[
\mathcal L(\bar x_{T},y^{\star})-\mathcal L(x^{\star},\bar y_{T})
\;\le\;
\frac{1}{2T}\Bigl(
\frac{\norm{x_{0}-x^{\star}}^{2}}{\tau}
+\frac{\norm{y_{0}-y^{\star}}^{2}}{\sigma}
\Bigr).
\tag{5}
\]

Consequently, $\bar x_{T}$ satisfies  

\[
f(\bar x_{T})+g(K\bar x_{T})-f(x^{\star})-g(Kx^{\star})\le O\!\left(\frac{1}{T}\right).
\]
\end{theorem}

\section*{Theoretical Properties}
\begin{itemize}
\item \textbf{Stability.} Condition \eqref{A2} guarantees that the joint Lyapunov functional
\[
\Phi_{k}:=\frac{1}{2\tau}\norm{x_{k}-x^{\star}}^{2}+\frac{1}{2\sigma}\norm{y_{k}-y^{\star}}^{2}
\]
is non‑increasing up to a controllable error term (see Lemma~\ref{lem:descent}).
\item \textbf{Hyper‑parameter sensitivity.} The bound \eqref{5} depends linearly on $\tau^{-1}$ and $\sigma^{-1}$. Choosing $\tau=\sigma=1/\|K\|$ (the largest admissible pair) minimises the constant factor.
\item \textbf{Comparison.} For pure gradient descent on $f$ (i.e. $g\equiv0$) the same $O(1/T)$ rate holds with stepsize $1/L$. PDHG recovers gradient descent when $K=I$ and $g\equiv0$, but additionally handles the extra term $g(Kx)$ without any inner loop.
\end{itemize}

\section*{Discussion}
The theorem shows that PDHG attains the optimal ergodic rate for convex optimisation without requiring strong convexity. The proof hinges on a simple three‑point identity and the cocoercivity of the gradient of $f$ (smoothness) together with the proximal optimality conditions for $g^{*}$. Extensions to the strongly convex case yield linear convergence, while stochastic variants preserve the $O(1/\sqrt{T})$ rate under variance bounds.

\newpage
\section*{Preliminaries (Definitions and Lemmas)}
\begin{lemma}[Three‑point identity for proximal operators]\label{lem:prox}
For any convex $h$ and any $u,v\in\mathbb R^{d}$,
\[
\langle u-\operatorname{prox}_{\lambda h}(u),\,\operatorname{prox}_{\lambda h}(u)-v\rangle
\ge \lambda\bigl(h(\operatorname{prox}_{\lambda h}(u))-h(v)\bigr).
\tag{L1}
\]
\end{lemma}
\begin{proof}
The optimality condition for $p:=\operatorname{prox}_{\lambda h}(u)$ reads  
$0\in \partial h(p)+\frac{1}{\lambda}(p-u)$, i.e.  
$\frac{1}{\lambda}(u-p)\in\partial h(p)$. By monotonicity of $\partial h$,
\[
\big\langle \tfrac{1}{\lambda}(u-p)-\xi,\;p-v\big\rangle\ge0,\qquad\forall \xi\in\partial h(v).
\]
Rearranging gives \eqref{L1}. ∎
\end{proof}

\begin{lemma}[Descent Lemma for $L$‑smooth $f$]\label{lem:smooth}
For all $x,x'$,
\[
f(x')\le f(x)+\langle\nabla f(x),x'-x\rangle+\frac{L}{2}\norm{x'-x}^{2}.
\tag{L2}
\]
\end{lemma}
\begin{proof}
Standard consequence of $L$‑smoothness; see Assumption~\ref{as:convex} (A1). ∎
\end{proof}

\begin{lemma}[Key descent inequality]\label{lem:descent}
Let $(x_{k},y_{k})$ be generated by \eqref{2a}–\eqref{2b}. Then for any saddle point $(x^{\star},y^{\star})$,
\[
\Phi_{k+1}\le \Phi_{k}
-\frac{1-\theta}{2\tau}\norm{x_{k+1}-x_{k}}^{2}
-\frac{1-\theta}{2\sigma}\norm{y_{k+1}-y_{k}}^{2}
+\mathcal L(x_{k+1},y^{\star})-\mathcal L(x^{\star},y_{k+1}),
\tag{L3}
\]
where $\Phi_{k}$ is defined in the previous section and $\theta=\tau\sigma\|K\|^{2}$.
\end{lemma}
\begin{proof}
\textbf{[Step 1]} Apply Lemma~\ref{lem:prox} to the dual update \eqref{2a} with $h=g^{*}$, $\lambda=\sigma$, $u=y_{k}+\sigma Kx_{k}$ and $v=y^{\star}$:
\[
\langle y_{k}+\sigma Kx_{k}-y_{k+1},\,y_{k+1}-y^{\star}\rangle
\ge\sigma\bigl(g^{*}(y_{k+1})-g^{*}(y^{\star})\bigr).
\tag{Step 1}
\]

\textbf{[Step 2]} Use the definition of the proximal operator for the primal step \eqref{2b} and Lemma~\ref{lem:prox} with $h=f$, $\lambda=\tau$, $u=x_{k}-\tau K^{\top}y_{k+1}$, $v=x^{\star}$:
\[
\langle x_{k}-\tau K^{\top}y_{k+1}-x_{k+1},\,x_{k+1}-x^{\star}\rangle
\ge\tau\bigl(f(x_{k+1})-f(x^{\star})\bigr).
\tag{Step 2}
\]

\textbf{[Step 3]} Expand the inner products:
\[
\begin{aligned}
\langle y_{k}-y_{k+1},y_{k+1}-y^{\star}\rangle
&+\sigma\langle Kx_{k},y_{k+1}-y^{\star}\rangle
\ge\sigma\bigl(g^{*}(y_{k+1})-g^{*}(y^{\star})\bigr),\\[1mm]
\langle x_{k}-x_{k+1},x_{k+1}-x^{\star}\rangle
&-\tau\langle K^{\top}y_{k+1},x_{k+1}-x^{\star}\rangle
\ge\tau\bigl(f(x_{k+1})-f(x^{\star})\bigr).
\end{aligned}
\tag{Step 3}
\]

\textbf{[Step 4]} Use the identity $\langle a-b,a-c\rangle=\frac12(\|a-b\|^{2}+\|a-c\|^{2}-\|b-c\|^{2})$ on both lines of \eqref{Step 3} to obtain
\[
\begin{aligned}
\frac{1}{2\sigma}\bigl(\|y_{k}-y^{\star}\|^{2}-\|y_{k+1}-y^{\star}\|^{2}
-\|y_{k+1}-y_{k}\|^{2}\bigr)
&+\langle Kx_{k},y_{k+1}-y^{\star}\rangle \\
&\ge g^{*}(y_{k+1})-g^{*}(y^{\star}),\\[1mm]
\frac{1}{2\tau}\bigl(\|x_{k}-x^{\star}\|^{2}-\|x_{k+1}-x^{\star}\|^{2}
-\|x_{k+1}-x_{k}\|^{2}\bigr)
&-\langle K^{\top}y_{k+1},x_{k+1}-x^{\star}\rangle \\
&\ge f(x_{k+1})-f(x^{\star}).
\end{aligned}
\tag{Step 4}
\]

\textbf{[Step 5]} Add the two inequalities of \eqref{Step 4}. The mixed terms combine:
\[
\langle Kx_{k},y_{k+1}\rangle-\langle Kx_{k},y^{\star}\rangle
-\langle K^{\top}y_{k+1},x_{k+1}\rangle+\langle K^{\top}y_{k+1},x^{\star}\rangle
=
\langle K(x_{k}-x_{k+1}),y_{k+1}\rangle
+\langle Kx_{k},y^{\star}\rangle-\langle Kx_{k+1},y^{\star}\rangle .
\tag{Step 5}
\]

\textbf{[Step 6]} Rearrange to the saddle‑point expression:
\[
\mathcal L(x_{k+1},y^{\star})-\mathcal L(x^{\star},y_{k+1})
=
f(x_{k+1})-f(x^{\star})+\langle Kx_{k+1},y^{\star}\rangle-g^{*}(y^{\star})
-\langle Kx^{\star},y_{k+1}\rangle+g^{*}(y_{k+1}).
\tag{Step 6}
\]

\textbf{[Step 7]} Substituting \eqref{Step 5}–\eqref{Step 6} into the sum of the two inequalities from \eqref{Step 4} yields
\[
\Phi_{k+1}\le\Phi_{k}
-\frac{1}{2\tau}\norm{x_{k+1}-x_{k}}^{2}
-\frac{1}{2\sigma}\norm{y_{k+1}-y_{k}}^{2}
+\mathcal L(x_{k+1},y^{\star})-\mathcal L(x^{\star},y_{k+1})
+\langle K(x_{k}-x_{k+1}),y_{k+1}-y^{\star}\rangle .
\tag{Step 7}
\]

\textbf{[Step 8]} Apply Cauchy–Schwarz and the definition $\theta=\tau\sigma\|K\|^{2}$:
\[
\langle K(x_{k}-x_{k+1}),y_{k+1}-y^{\star}\rangle
\le \|K\|\norm{x_{k+1}-x_{k}}\norm{y_{k+1}-y^{\star}}
\le \frac{\theta}{2}\Bigl(\frac{1}{\tau}\norm{x_{k+1}-x_{k}}^{2}
+\frac{1}{\sigma}\norm{y_{k+1}-y^{\star}}^{2}\Bigr).
\tag{Step 8}
\]

\textbf{[Step 9]} Insert the bound \eqref{Step 8} into \eqref{Step 7}. Collecting the coefficients of the norm terms yields exactly inequality \eqref{L3}. ∎
\end{proof}

\section*{Main Proof (Step‑by‑Step Derivation)}
We now prove Theorem~\ref{thm:ergodic}.

\textbf{[Step 1]} Sum Lemma~\ref{lem:descent} from $k=0$ to $T-1$:
\[
\sum_{k=0}^{T-1}\Bigl[\Phi_{k+1}-\Phi_{k}\Bigr]
\le
-\frac{1-\theta}{2\tau}\sum_{k=0}^{T-1}\norm{x_{k+1}-x_{k}}^{2}
-\frac{1-\theta}{2\sigma}\sum_{k=0}^{T-1}\norm{y_{k+1}-y_{k}}^{2}
+\sum_{k=0}^{T-1}\bigl[\mathcal L(x_{k+1},y^{\star})-\mathcal L(x^{\star},y_{k+1})\bigr].
\tag{S1}
\]

\textbf{[Step 2]} The left‑hand side telescopes:
\[
\Phi_{T}-\Phi_{0}\le
-\frac{1-\theta}{2\tau}\sum_{k=0}^{T-1}\norm{x_{k+1}-x_{k}}^{2}
-\frac{1-\theta}{2\sigma}\sum_{k=0}^{T-1}\norm{y_{k+1}-y_{k}}^{2}
+\sum_{k=0}^{T-1}\bigl[\mathcal L(x_{k+1},y^{\star})-\mathcal L(x^{\star},y_{k+1})\bigr].
\tag{S2}
\]

\textbf{[Step 3]} Drop the non‑positive norm sums (they are $\le0$) to obtain an upper bound:
\[
\sum_{k=0}^{T-1}\bigl[\mathcal L(x_{k+1},y^{\star})-\mathcal L(x^{\star},y_{k+1})\bigr]
\le \Phi_{0}-\Phi_{T}
\le \Phi_{0}.
\tag{S3}
\]

\textbf{[Step 4]} By definition of $\Phi_{0}$,
\[
\Phi_{0}= \frac{1}{2\tau}\norm{x_{0}-x^{\star}}^{2}
+\frac{1}{2\sigma}\norm{y_{0}-y^{\star}}^{2}.
\tag{S4}
\]

\textbf{[Step 5]} Divide inequality \eqref{S3} by $T$ and use convexity of $\mathcal L$ in its first argument and concavity in its second argument:
\[
\mathcal L\!\bigl(\bar x_{T},y^{\star}\bigr)-\mathcal L\!\bigl(x^{\star},\bar y_{T}\bigr)
\le \frac{1}{T}\sum_{k=0}^{T-1}\bigl[\mathcal L(x_{k+1},y^{\star})-\mathcal L(x^{\star},y_{k+1})\bigr]
\le \frac{\Phi_{0}}{T}.
\tag{S5}
\]

\textbf{[Step 6]} Substituting \eqref{S4} gives the desired bound \eqref{5}:
\[
\mathcal L(\bar x_{T},y^{\star})-\mathcal L(x^{\star},\bar y_{T})
\le \frac{1}{2T}\Bigl(
\frac{\norm{x_{0}-x^{\star}}^{2}}{\tau}
+\frac{\norm{y_{0}-y^{\star}}^{2}}{\sigma}
\Bigr).
\tag{S6}
\]

\textbf{[Step 7]} The saddle‑point definition $\mathcal L(x^{\star},y^{\star})\le\mathcal L(\bar x_{T},y^{\star})$ and $\mathcal L(x^{\star},\bar y_{T})\le\mathcal L(x^{\star},y^{\star})$ imply
\[
f(\bar x_{T})+g(K\bar x_{T})-f(x^{\star})-g(Kx^{\star})
\le \mathcal L(\bar x_{T},y^{\star})-\mathcal L(x^{\star},\bar y_{T}),
\]
hence the $O(1/T)$ rate for the primal objective follows directly from \eqref{S6}. ∎

\section*{Rate Derivation (With Constants)}
From \eqref{S6} we read off the explicit constant
\[
C:=\frac12\Bigl(\frac{\norm{x_{0}-x^{\star}}^{2}}{\tau}
+\frac{\norm{y_{0}-y^{\star}}^{2}}{\sigma}\Bigr),
\qquad
\text{so that}\quad
\mathcal L(\bar x_{T},y^{\star})-\mathcal L(x^{\star},\bar y_{T})\le \frac{C}{T}.
\]
Choosing the largest admissible stepsizes, i.e. $\tau=\sigma=1/\|K\|$ (which saturates \eqref{A2} with $\theta=1/2$), minimises $C$ and yields the tightest bound.

\section*{Special Cases}
\begin{enumerate}
\item \textbf{No dual term} ($g\equiv0$). Then $y_{k}=0$ for all $k$, and \eqref{2b} reduces to the standard gradient descent step $x_{k+1}=x_{k}-\tau\nabla f(x_{k})$. The theorem recovers the classic $O(1/T)$ ergodic rate for smooth convex $f$.
\item \textbf{Strongly convex $f$}. If $f$ is $\mu$‑strongly convex, the Lyapunov functional can be strengthened to obtain linear convergence:
\[
\Phi_{k}\le (1-\rho)^{k}\Phi_{0},
\qquad\rho:=\frac{\mu\tau}{1+\mu\tau}.
\]
\item \textbf{Nonsmooth $g$}. The proof above uses only the proximal optimality condition of $g^{*}$, therefore the same $O(1/T)$ bound holds when $g$ is merely convex (e.g. $g=\|\cdot\|_{1}$).
\end{enumerate}

\end{document}